\chapter{A Causal Layer: Does Financing Cause Growth?}
\label{ch:causal}

Chapter~\ref{ch:xgb} predicts. Managers ask a different question: \emph{if} a company
raises a round, does growth follow --- or do growing companies simply raise? The fitted
$\hat m(x)$ cannot answer this, because the coefficientless dependence it learns on
deal-history features mixes the effect of financing with the selection of good firms into
financing. This chapter builds the minimal causal layer the European regime supports. The
regime is favourable for exactly one reason: the statutory fundamentals panel gives us a
\emph{rich, audited, pre-treatment} covariate history --- close to the dossier an
investor actually screens on --- so adjustment has a fighting chance. Every estimate
ships with cluster-bootstrap uncertainty, an E-value, and two falsification tests, and
the whole pipeline is validated on the synthetic testbed of Chapter~\ref{ch:selection},
where the true effect is known by construction.

\section{Estimands and assumptions}
\label{sec:causal-estimands}

Treatment: $A_{i,y} = \ind{\text{firm } i \text{ closed a PE/VC equity round in year }
y}$, from the deals table. Outcome: $Y_{i,y}$, the $h=2$ winsorized EBITDA arc-growth of
Chapter~\ref{ch:success}, measured over $(y, y+h]$ so it is strictly post-treatment.
Each unit carries potential outcomes $Y_{i,y}(1)$ and $Y_{i,y}(0)$: growth with and
without a round in year $y$.

\begin{assumption}[SUTVA]
\label{ass:causal-sutva}
(i)~No interference: $Y_{i,y}(a)$ does not depend on other firms' treatments;
(ii)~single version of treatment: $A_{i,y}=1$ denotes one well-defined intervention
\cite{imbens2015}.
\end{assumption}

\begin{assumption}[Consistency]
\label{ass:causal-consistency}
$Y_{i,y} = A_{i,y} Y_{i,y}(1) + (1 - A_{i,y})\, Y_{i,y}(0)$ \cite{hernan2020}.
\end{assumption}

Both parts of Assumption~\ref{ass:causal-sutva} are approximations here: a round can
hurt a direct competitor's growth (interference; we flag narrow niches), and rounds
differ in size and investor (versions; we deliberately coarsen to ``raised
institutional equity'' and treat dose--response as an extension, not the estimand).

The estimands, in decreasing ambition of the conditioning:
\begin{equation}
\tau_{\mathrm{ATE}} = \E[Y(1) - Y(0)], \quad
\tau_{\mathrm{ATT}} = \E[Y(1) - Y(0) \mid A = 1], \quad
\tau(x) = \E[Y(1) - Y(0) \mid X = x].
\label{eq:causal-estimands}
\end{equation}
For the portfolio question (``did our financings work?'') the ATT is the object; for
targeting (``where would a round move the needle?'') the CATE $\tau(x)$; the ATE is the
reporting default because it is the easiest to estimate stably. Identification requires:

\begin{assumption}[Conditional exchangeability and positivity]
\label{ass:causal-ignorability}
There exists an observed pre-treatment set $Z$ with
$\bigl(Y(0), Y(1)\bigr) \indep A \mid Z$ and
$0 < e(z) = \Prob(A = 1 \mid Z = z) < 1$ for all $z$ in the support
\cite{imbens2015,hernan2020}.
\end{assumption}

Assumption~\ref{ass:causal-ignorability} is not a modelling choice; it is a claim about
the world that the next section argues from a graph, Section~\ref{sec:causal-sens}
stress-tests, and Section~\ref{sec:causal-synth} verifies where the truth is known.

\section{The causal graph}
\label{sec:causal-dag}

Five variables. $U$: latent firm quality (management, product, moat) --- unobserved.
$X$: the observed pre-treatment fundamentals history (levels, growth path, margins,
leverage, payroll over the three years up to $y$), plus sector and age. $M$: the macro
state as of year $y$ (rates, CPI YoY, unemployment, as-of vintages). $A$: the treatment.
$Y$: the outcome. The edges: $U \to A$ (investors select quality), $U \to Y$ (quality
drives growth), $U \to X$ (quality leaves tracks in the filings), $X \to A$ (investors
screen on the observable dossier), $X \to Y$ (fundamentals persist), $M \to A$ (funding
cycles), $M \to Y$ (macro drives growth), and the causal edge of interest $A \to Y$.
Table~\ref{tab:causal-paths} enumerates every $A$--$Y$ path.

\begin{table}[htbp]
\centering
\caption{Paths between $A$ and $Y$ in the working graph.}
\label{tab:causal-paths}
\begin{tabular}{@{}llll@{}}
\toprule
Path & Type & Status after conditioning on $(X, M)$ \\
\midrule
$A \to Y$ & causal & kept (the estimand) \\
$A \leftarrow X \to Y$ & backdoor & blocked by $X$ \\
$A \leftarrow M \to Y$ & backdoor & blocked by $M$ (sector-year cells) \\
$A \leftarrow U \to Y$ & backdoor & \emph{not} blockable directly; attenuated \\
 & & via the proxy path $U \to X$ \\
$A \leftarrow U \to X \to Y$ & backdoor & blocked by $X$ \\
\bottomrule
\end{tabular}
\end{table}

The whole argument of the EU track sits in the fourth row. $U$ is unobserved, so the
path $A \leftarrow U \to Y$ cannot be blocked exactly. But investors do not observe $U$
either: they observe the dossier, which is essentially $X$ --- audited statements,
growth trajectory, margins, leverage. To the extent that the selection $U \to A$ flows
\emph{through} what is observable about quality, conditioning on a rich $X$ intercepts
most of it, and the residual confounding is the part of quality that moves investors
without leaving any track in three years of filings. In the European regime that
residual is plausibly small and, crucially, \emph{boundable}
(Section~\ref{sec:causal-sens}). In the US regime there is no fundamentals panel, $X$
collapses to deal marks and coarse firmographics, and this argument dies --- which is
why Chapters~\ref{ch:us} and onward change the estimand rather than pretend the backdoor
is closed.

\section{Identification: the backdoor criterion}
\label{sec:causal-backdoor}

\begin{definition}[Backdoor criterion \cite{pearl1995}]
\label{def:causal-backdoor}
A set of observed variables $Z$ satisfies the backdoor criterion relative to $(A, Y)$
if (i)~no element of $Z$ is a descendant of $A$, and (ii)~$Z$ blocks every path between
$A$ and $Y$ that contains an arrow into $A$.
\end{definition}

\begin{theorem}[Adjustment formula \cite{pearl1995}]
\label{thm:causal-adjustment}
If $Z$ satisfies the backdoor criterion relative to $(A, Y)$, then
\begin{equation}
\E[Y \mid \Do(A = a)] = \E_Z\bigl[\, \E[Y \mid A = a, Z] \,\bigr],
\label{eq:causal-adjustment}
\end{equation}
the outer expectation over the marginal law of $Z$, and consequently
$\tau_{\mathrm{ATE}} = \E_Z[\E[Y \mid A=1, Z] - \E[Y \mid A=0, Z]]$.
\end{theorem}

In the working graph, $Z = (X, M)$ satisfies the criterion \emph{up to} the residual
$U$-path discussed above: formally we proceed under
Assumption~\ref{ass:causal-ignorability} with $Z = (X$ history over 3 years, sector-year
cell$)$ --- the cells absorb $M$ and sector-level funding cycles --- and quantify the
violation rather than assume it away.

\begin{warning}[Bad controls]
\label{warn:causal-badcontrols}
More conditioning is not more safety. Three failure modes, all present in this dataset.
(i)~\emph{Mediators}: post-treatment employee counts sit on the path
$A \to \text{hiring} \to Y$; conditioning on them subtracts part of the effect being
estimated. (ii)~\emph{Colliders/descendants}: a subsequent valuation mark is caused
both by the round ($A$) and by quality ($U$); conditioning on it opens the path
$A \to \text{mark} \leftarrow U \to Y$ and manufactures spurious association.
(iii)~\emph{M-bias}: even a \emph{pre-treatment} variable can hurt --- if some early
covariate is a common effect of one latent cause that drives selection and another
latent cause that drives growth, then conditioning on it opens a previously closed path
between those latents. The rule is graphical, not chronological: the adjustment set is
chosen by Definition~\ref{def:causal-backdoor}, never by predictive power, and nothing
dated after the round enters $Z$.
\end{warning}

\section{Estimators, in increasing robustness}
\label{sec:causal-estimators}

Throughout, nuisance functions are fitted with the boosted-tree machinery of
Chapter~\ref{ch:xgb} (same features, same era discipline), written
$\mu_a(z) = \E[Y \mid A = a, Z = z]$ and $e(z) = \Prob(A = 1 \mid Z = z)$.

\paragraph{(1) Outcome regression (g-formula).}
Fit $\hat\mu_1, \hat\mu_0$ and average:
$\hat\tau_{\mathrm{OR}} = n^{-1} \sum_{i,y} \{\hat\mu_1(Z_{i,y}) -
\hat\mu_0(Z_{i,y})\}$. Two flavours: the T-learner (separate models per arm; honest but
noisy in the small treated arm) and the S-learner ($A$ as one feature among $p$; stable
but regularization can shrink the $A$-splits, biasing $\hat\tau$ toward zero). Both
inherit every extrapolation weakness of trees precisely in low-overlap regions, where
$\hat\mu_{1}$ is evaluated at control-like $z$ and vice versa.

\paragraph{(2) Inverse propensity weighting.}
Fit $\hat e$, then the Hajek form
\[
\hat\tau_{\mathrm{IPW}}
= \frac{\sum_i A_i Y_i w_i^1}{\sum_i A_i w_i^1}
- \frac{\sum_i (1-A_i) Y_i w_i^0}{\sum_i (1-A_i) w_i^0},
\qquad
w^1 = \frac{1}{\hat e(Z)}, \quad w^0 = \frac{1}{1 - \hat e(Z)}.
\]
Mandatory diagnostics before any estimate is read:
the histogram of $\hat e$ by treatment arm (overlap plot); trimming to
$\hat e \in [0.02, 0.98]$ with the trimmed fraction reported --- if it exceeds a few
percent, the estimand has silently changed to a trimmed-population effect and we say
so; the effective sample size $\mathrm{ESS} = (\sum w)^2 / \sum w^2$ per arm. IPW uses
the treatment model only, so it fails independently of (1) --- but its variance explodes
near the trimming boundary.

\paragraph{(3) AIPW with cross-fitting (the default).}
The doubly robust score combines both nuisances \cite{chernozhukov2018}:
\begin{equation}
\psi(O; \mu, e) = \mu_1(Z) - \mu_0(Z)
+ \frac{A\,(Y - \mu_1(Z))}{e(Z)}
- \frac{(1-A)\,(Y - \mu_0(Z))}{1 - e(Z)},
\qquad
\hat\tau_{\mathrm{AIPW}} = \frac{1}{n} \sum \psi_i .
\label{eq:causal-aipw}
\end{equation}

\begin{proposition}[Double robustness]
\label{prop:causal-dr}
Under Assumptions~\ref{ass:causal-consistency}--\ref{ass:causal-ignorability},
$\E[\psi(O; \mu, e)] = \tau_{\mathrm{ATE}}$ if \emph{either} $\mu = (\mu_1^*, \mu_0^*)$
is the true outcome function \emph{or} $e = e^*$ is the true propensity (with the other
nuisance arbitrary, subject to positivity).
\end{proposition}

\begin{proof}
Consider the $a=1$ half, $\psi_1 = \mu_1(Z) + A(Y - \mu_1(Z))/e(Z)$. If
$\mu_1 = \mu_1^*$: $\E[A(Y - \mu_1^*(Z)) \mid Z] = e^*(Z)\,\E[Y - \mu_1^*(Z) \mid A=1,
Z] = 0$, so $\E[\psi_1] = \E[\mu_1^*(Z)] = \E[Y(1)]$ by
Assumption~\ref{ass:causal-ignorability}. If $e = e^*$:
$\E[\psi_1 \mid Z] = \mu_1(Z) + \{e^*(Z)\,\E[Y \mid A=1,Z] - e^*(Z)\mu_1(Z)\}/e^*(Z)
= \E[Y \mid A=1, Z] = \E[Y(1) \mid Z]$, and the arbitrary $\mu_1$ cancels. The $a=0$
half is symmetric; subtract.
\end{proof}

\begin{remark}[Rate condition]
\label{rem:causal-rates}
Beyond consistency, valid $\sqrt{n}$ inference informally requires the \emph{product}
of nuisance errors to vanish fast:
$\|\hat\mu_a - \mu_a^*\|_{2} \cdot \|\hat e - e^*\|_{2} = o_P(n^{-1/2})$
\cite{chernozhukov2018}. Each nuisance may converge slowly (as boosted trees do) as
long as the product is fast --- this, plus the removal of own-observation overfitting
bias by cross-fitting, is why (3) is the default and not merely a hedge.
\end{remark}

\paragraph{Algorithm (cross-fitted AIPW).}
\begin{enumerate}
\item Partition \emph{firms} into $K = 5$ folds; a firm's rows never straddle folds
(within-firm correlation would otherwise leak nuisance fits into own scores).
\item For each fold $k$: fit $\hat\mu_1^{(-k)}, \hat\mu_0^{(-k)}, \hat e^{(-k)}$ on the
other $K-1$ folds, with the era-aware tuning of Chapter~\ref{ch:xgb} frozen once,
outside the loop.
\item Trim: drop rows of fold $k$ with $\hat e^{(-k)} \notin [0.02, 0.98]$; record the
trimmed fraction.
\item Compute $\psi_{i,y}$ for the surviving rows of fold $k$ via
\eqref{eq:causal-aipw} using only out-of-fold nuisances.
\item Pool: $\hat\tau = $ mean of all $\psi_{i,y}$; variance by
Section~\ref{sec:causal-inference}; report alongside the overlap plot and trimmed
fraction.
\end{enumerate}

\section{Inference under firm clustering}
\label{sec:causal-inference}

Rows recur within firms across years, and the scores $\psi_{i,y}$ of a firm are
correlated through persistent firm effects. An i.i.d.\ bootstrap (or the naive variance
$n^{-2}\sum(\psi_{i,y} - \hat\tau)^2$) treats $n$ rows as $n$ independent draws and
understates the variance --- the effective sample size is nearer the number of firms.
Two consistent routes; we run both and report the wider.

\paragraph{(a) Cluster-robust influence-function variance.}
Let $S_i = \sum_{y} (\psi_{i,y} - \hat\tau)$ sum a firm's centred scores. Then
\begin{equation}
\widehat{\Var}(\hat\tau) = \frac{1}{n^2} \sum_{i \in \mathcal{F}} S_i^{\,2},
\label{eq:causal-clustervar}
\end{equation}
which is the sandwich variance with firms as clusters; normal intervals
$\hat\tau \pm z_{1-\alpha/2}\,\widehat{\Var}^{1/2}$.

\paragraph{(b) Cluster bootstrap \cite{efron1994}.}
\begin{enumerate}
\item For $b = 1, \dots, B$ (200--500): resample firms with replacement; a drawn firm
contributes all its rows.
\item Re-run the \emph{entire} cross-fitting algorithm on the replicate --- nuisance
refits included, tuning still frozen --- to obtain $\hat\tau_b$.
\item Report the percentile interval of $\{\hat\tau_b\}$.
\end{enumerate}
The bootstrap additionally propagates nuisance-estimation noise that
\eqref{eq:causal-clustervar} treats as fixed; disagreement between (a) and (b) beyond
$\sim$20\% of interval width is itself a diagnostic of nuisance instability.

\section{Sensitivity and falsification}
\label{sec:causal-sens}

Assumption~\ref{ass:causal-ignorability} cannot be verified from the data it licenses.
It can be stress-tested. Every reported effect must clear the checklist of
Table~\ref{tab:causal-checklist}; an estimate that fails any row does not ship.

\paragraph{E-values.}
The E-value \cite{vanderweele2017} is the minimum strength of association, on the risk
ratio scale, that an unmeasured confounder would need with \emph{both} treatment and
outcome, above and beyond $Z$, to explain away the observed association. For an
observed risk ratio $\mathrm{RR} \ge 1$,
\begin{equation}
\mathrm{E} = \mathrm{RR} + \sqrt{\mathrm{RR}\,(\mathrm{RR} - 1)} .
\label{eq:causal-evalue}
\end{equation}
Our outcome is continuous, so we standardize: with effect $d = \hat\tau / \hat\sigma_Y$
(cluster-bootstrap CI carried through), the conventional conversion
$\mathrm{RR} \approx \exp(0.91\, d)$ feeds \eqref{eq:causal-evalue}, and the E-value of
the CI limit closer to the null is reported alongside that of the point estimate. An
E-value of 1.8 reads: a latent-quality confounder must nearly double both the odds of
raising and of high growth, conditional on the full fundamentals history, to explain
the estimate away --- a statement a domain expert can argue with, which is the point.

\paragraph{Placebo treatment.}
Define $A^{\mathrm{fut}}_{i,y} = \ind{\text{round in year } y + 3}$ and estimate its
``effect'' on $Y_{i,y}$ (which is fully realized before the placebo round) with the
identical AIPW pipeline. Under Assumption~\ref{ass:causal-ignorability} the estimate
must be null: a future round cannot cause past growth, so any nonzero adjusted estimate
measures exactly the selection-on-trends that $Z$ failed to absorb.

\paragraph{Negative-control outcome.}
Symmetrically, estimate the ``effect'' of $A_{i,y}$ on \emph{pre-treatment} growth
(years $y-2$ to $y$, excluded from $Z$ for this test only). Nonzero means the treated
were already on a different trajectory in a way $Z$ does not capture.

\begin{table}[htbp]
\centering
\caption{Mandatory pre-shipping checklist for any causal estimate.}
\label{tab:causal-checklist}
\small
\begin{tabular}{@{}lll@{}}
\toprule
Check & Pass criterion & On failure \\
\midrule
Overlap plot & common support, trimmed & redefine estimand (ATT, \\
 & fraction $< 5\%$ & trimmed population) \\
Placebo (future round) & CI covers 0 & enrich $Z$ (longer history, \\
 & & trend features); else do not ship \\
Negative control (pre-growth) & CI covers 0 & same as placebo \\
E-value & $>$ plausible latent-quality & report as bounded, not \\
 & association given rich $X$ & as an effect \\
Variance agreement & IF vs bootstrap within 20\% & investigate nuisance \\
 & & instability \\
Synthetic recovery & Table~\ref{tab:causal-synth} ranking & pipeline bug; fix before \\
 & reproduced & touching real data \\
\bottomrule
\end{tabular}
\end{table}

\section{Synthetic validation protocol}
\label{sec:causal-synth}

The testbed of Chapter~\ref{ch:selection} generates a world with a known latent quality
$q_i$ driving both funding intensity and fundamentals, so the true $\tau$ is known by
construction and confounding by quality is present by design. The causal layer is
accepted only if it passes the following protocol there.

\begin{enumerate}
\item \emph{Generate} an EU-style panel: fundamentals as noisy, lagged functions of
$q_i$ and macro; treatment assignment with probability increasing in $q_i$ and in the
observable dossier; outcome with a known additive effect $\tau^* $ of treatment.
\item \emph{Estimate} with four pipelines: naive difference in means; g-formula
(T-learner); IPW; cross-fitted AIPW --- each with firm-cluster bootstrap CIs and
E-values.
\item \emph{Verify} the qualitative ranking of Table~\ref{tab:causal-synth}: the naive
contrast overstates $\tau^*$ badly (quality confounding); all adjusted estimators land
near $\tau^*$ when the generated $X$ is rich; ablating $X$ columns (hiding the dossier,
i.e.\ simulating the US regime) re-opens the gap, with AIPW degrading most gracefully.
\item \emph{Report} empirical CI coverage of $\tau^*$ across Monte Carlo replications
and check that the naive estimator's E-value correctly flags fragility while the
adjusted estimators' E-values are informative about the injected residual confounding.
\end{enumerate}

\begin{table}[htbp]
\centering
\caption{Expected qualitative behaviour on the synthetic testbed
(Chapter~\ref{ch:selection}). ``$\approx$'' means CI covers $\tau^*$ at nominal rate.}
\label{tab:causal-synth}
\small
\begin{tabular}{@{}lll@{}}
\toprule
Estimator & Rich $X$ (EU-style) & Impoverished $X$ (US-style) \\
\midrule
Naive difference & $\gg \tau^*$ & $\gg \tau^*$ \\
g-formula (XGB) & $\approx \tau^*$ & biased toward naive \\
IPW (trimmed) & $\approx \tau^*$, wider CI & biased, unstable \\
AIPW cross-fitted & $\approx \tau^*$, correct coverage & residual bias, flagged \\
 & & by placebo and E-value \\
\bottomrule
\end{tabular}
\end{table}

This protocol is what makes the chapter testable rather than decorative: the code that
produces Table~\ref{tab:causal-synth} is byte-identical to the code run on the real
panel, so a pipeline that cannot recover a known effect is caught before it is believed.

\section{Limits}
\label{sec:causal-limits}

Two limits are structural and stated plainly. First, \emph{selection on unobservables
is not solved} --- it is attenuated by a rich $X$, bounded by E-values, and probed by
placebo and negative-control tests, but a confounder invisible to three years of
audited filings and strong enough to survive the checklist remains logically possible.
Our claim is calibrated accordingly: adjusted estimates with reported sensitivity, not
oracle effects. Second, \emph{external validity across macro regimes}: $M$ enters both
selection and outcome, all estimates average over the funding cycles in sample, and a
CATE-by-era analysis is data-hungry precisely because treated firms are scarce within
any single era. An effect estimated across 2015--2021 need not transport to a
tight-money regime, for the same no-extrapolation reasons as in
Chapter~\ref{ch:xgb}. The US regime sharpens both limits at once: with no fundamentals
panel there is no rich $X$, and the honest response is to change the estimand and the
observation model --- the programme of Chapters~\ref{ch:us} and~\ref{ch:selection}.
